\documentclass[oneside,a4paper,14pt]{extarticle} %размер шрифта 14
\usepackage[T1,T2A]{fontenc}
\usepackage[
a4paper,
letterpaper,
top=2cm,
bottom=2cm,
left=2.5cm,
right=1.5cm
]{geometry}
\usepackage[utf8]{inputenc} %кодировка текста
\usepackage[russian]{babel} %поддержка русского языка
\usepackage{textcomp} %текстовые символы
\usepackage{indentfirst} %корректировка отступов
\usepackage{listings} %для кода !!
\usepackage{float} % для "впихивания" изображений в странницу
\usepackage{xcolor}
\usepackage{enumitem} %для нумерации буквами
\usepackage{graphicx} %работа с изображениям
\usepackage{mwe} % for blindtext and example-image-a in example
\usepackage{wrapfig}
\usepackage{caption}
\usepackage{amsmath}  % для формул и символов
\usepackage{amsfonts}
\usepackage{amsthm}
\usepackage{graphicx}
\usepackage[all]{xy}
\usepackage[breaklinks]{hyperref}
\usepackage{xurl}

%%размеркегляузаголоковразделов
\usepackage{titlesec}
\titleformat{\section}
{\normalsize\bfseries}
{\thesection} {1em}{}
\titleformat{\subsection}
{\normalsize\bfseries}
{\thesubsection} {1em}{}
\titleformat{\subsubsection}
{\normalsize\bfseries}
{\thesubsection} {1em}{}

\renewcommand\baselinestretch{1.33}\normalsize % межстрочный интервал
\setlength{\parindent}{1.25cm}
\usepackage{indentfirst}

\lstset{
	language=C,
	basicstyle=\ttfamily\small,
	keywordstyle=\color{blue},
	commentstyle=\color{green!60!black},
	stringstyle=\color{red},
	numbers=left,
	numberstyle=\tiny\color{gray},
	frame=single,
	breaklines=true,
	showstringspaces=false,
	tabsize=2,
	inputencoding=utf8,
	extendedchars=true
}

\begin{document}
	\newpage\thispagestyle{empty}
	\begin{center}
		МИНИСТЕРСТВО НАУКИ И ВЫСШЕГО ОБРАЗОВАНИЯ\\
		РОССИЙСКОЙ ФЕДЕРАЦИИ
		ФЕДЕРАЛЬНОЕ ГОСУДАРСТВЕННОЕ БЮДЖЕТНОЕ\\
		ОБРАЗОВАТЕЛЬНОЕ
		УЧРЕЖДЕНИЕ ВЫСШЕГО ОБРАЗОВАНИЯ\\
		«ВЯТСКИЙ ГОСУДАРСТВЕННЫЙ УНИВЕРСИТЕТ»\\
		Институт математики и информационных систем\\
		Факультет автоматики и вычислительной техники\\
		Кафедра электронных вычислительных машин
	\end{center}
	\vspace{20mm}
	
	\begin{center}
		Отчёт по лабораторной работе №6\\
		по дисциплине\\
		<<Алгоритмы и структуры данных>>\\
	\end{center}
	\vspace{48mm}
	
	\noindent
	Выполнил студент гр. ИВТб-1303-06-00 \hfill \rule[-0,5mm]{30mm}{0.15mm}\,/Шубин В.Ю./
	\newline
	Проверил заведующий кафедрой ЭВМ \hfill \rule[-0,5mm]{30mm}{0.15mm}\,/Долженкова М.Л./
	
	\vfill
	\begin{center}
		Киров\\
		2026
	\end{center}
	
	\newpage
	
	\section*{Цель работы:}
	Отточить навыки решения задач в условиях ограниченного времени.
	
	\section*{Вводные}
	Для выполнения данной лабораторной работы был предоставлен \\ Яндекс.Контест
	\url{https://official.contest.yandex.ru/contest/94203}
	
	\section*{Описание процесса решения}
	Был запущен контест, который длится 2 часа.
	Я приступил к решению первой задачи: A. Запрещенные биграммамы
	(\url{https://official.contest.yandex.ru/contest/94203/problems/A/}).
	Судя по условию, задача сводится к подсчёту количества строк длины N, в которых отсутствуют два одинаковых 
	символа, идущих подряд. По сути, множители просто чередуются между 1 и 3 (например, последовательность 1×3×1×3... 
	при старте с гласной, либо 3×1×3×1... в противном случае).
	Таким образом, для чётного N результат вычисляется как $2 \cdot 3^{n \bmod 2}$, а для нечётного 
	N — как $4 \cdot 3^{n \bmod 2}$. При этом итоговое значение необходимо вывести по модулю $10^9 + 7$. \\
	
	\lstinputlisting[language=Python, caption={Решение задачи на Python}]{1.py}
	Решение данной задачи заняло примерно 30 минут.
	
	\newpage
	
	Далее я приступил к решению задачи B. RL-агент стремится к хаосу
	(\url{https://official.contest.yandex.ru/contest/94203/problems/B/}). 
	В данной задаче нужно определить, есть ли в массиве разница между i-ым элементом и индексом i
	больше k, если да то вывести 1 иначе 0. Чтобы не проверять весь массив после каждой  перестановки
	я решил проверять каждую перестановку, сделала ли перестановка массив <<хаотичным>>. \\
	
	\lstinputlisting[caption={Решение задачи на C}]{2.c}
	Решение данной задачи заняло тоже примерно 30 минут. \\
	
	После решения первых двух задач остался 1 час соревнования. Я прочитал задание
	C. Рекомендательная система доставки
	(\url{https://official.contest.yandex.ru/contest/94203/problems/C/}).
	Поскольку общее число артикулов не превышает 100 тысяч, я сразу выделил под них фиксированный массив. 
	Каждая ячейка хранит пару целых чисел: стоимость доставки и номер склада.
	Начальное заполнение такое: цену доставки ставлю на максимально возможное значение, а номер склада — в 0 
	(так как реальная нумерация складов начинается с единицы).
	При чтении данных по складам для каждого артикула сверяюсь с текущей записью в массиве. Обновляю её, если 
	новая стоимость доставки строго меньше сохранённой, либо если цены совпадают, но у нового склада номер меньше.
	Затем перехожу к обработке K запросов. В каждом запросе приходят артикул и базовая цена товара. Проверяю массив: 
	если номер склада для этого артикула остался равным 0, значит, товар отсутствует — вывожу `-1 -1`. 
	Иначе считаю итоговую сумму (цена товара + сохранённая стоимость доставки) и вывожу результат. 
	\newpage
	\lstinputlisting[caption={Решение задачи на C}]{3.c}
	Решение третьей задачи заняло почти час. \\
	
	У меня уже не оставалось времени, но взглянув на оставшиеся задачи я понял, что такое мне
	точно не по силам
	
	\section*{Выводы}
	Подводя итоги можно сказать, что для успешного решения задач в контесте нужно уметь быстро 
	анализировать саму суть задачи, отбрасывая лишнее описание. Как показала практика, не
	всегда удаётся быстро понять смысл задания и придумать алгоритм решения. 
\end{document}